\documentclass[english, parskip=full]{scrartcl}
\usepackage[utf8]{inputenc}  % Kodierung der Datei
\usepackage[english]{babel}  % Sprache auf Deutsch 
\usepackage{color}
\usepackage{graphicx, gensymb}
\usepackage{amsfonts, cancel, amssymb, slashed}
\usepackage{amsmath, amsthm, array, esvect, esint}
\usepackage{booktabs, multirow}  % schönere Tabellen
\usepackage{caption}  % Anpassung der Captions
\usepackage{scrlayer-scrpage}    % Manipulation des Seitenstils
\pagestyle{scrheadings}  % Stil für die Seite setzen
\automark{section}  % Abschnittsnamen als Seitenbeschriftung 
\setheadsepline{.5pt}    % Kopzeile durch Linie abtrennen
\usepackage[hidelinks]{hyperref}  % Links und weitere PDF-Features
\usepackage{typearea, tikz, pgffor, verbatim}
\usetikzlibrary{calc, quotes}
\usepackage[cache=false, outputdir=build]{minted}
\usepackage{dsfont}


\definecolor{bg}{rgb}{0.9,0.9,0.9}
\newcommand\numberthis{\addtocounter{equation}{1}\tag{\theequation}}
\newcommand{\bra}{}
\newcommand{\kom}{}
\newcommand{\brak}{}
\newcommand{\braket}{}
\newcommand{\intx}{}
\newcommand{\intr}{}
\newcommand{\kla}{}
\newcommand{\klam}{}
\newcommand{\klammer}{}
\newcommand{\oper}{}
\newcommand{\tecxt}{}
\newcommand{\Bra}{}
\newcommand{\Ket}{}
\def\I{\text{i}}
\def\E{\text{e}}

\newcommand*{\titel}{07 Ladder operators}
\newcommand*{\autor}{Joachim Schwardt}

\hypersetup{pdfauthor={\autor}, pdftitle={\titel}}

\subject{Quantum theory 2}
\title{\titel}
\author{\autor}
\date{\today}

\begin{document}

%\begin{titlepage}
\maketitle  % Titel setzen
%\tableofcontents  % Inhaltsverzeichnis setzen
%\end{titlepage}

\section{Questions}
\section{Creation and annihilation operators}
We have $a_n^\dagger\Ket | {n_j} \rangle^\pm=\sqrt{N+1}\Ket | {n n_j} \rangle^\pm$ for $j=1,\dots,N$.
Hence
\begin{align*}
a_5^\dagger a_3^\dagger a_4^\dagger a_2^\dagger a_7^\dagger &=\sqrt{5!}\Ket | {5,3,4,2,7} \rangle^\pm
\end{align*}
We have $a_n\Ket | {n_j} \rangle^\pm=\frac{1}{\sqrt{N}}\sum_j\delta_{nn_j}(\pm 1)^{j+1}\Ket | {n_{j-1}n_{j+1}} \rangle^\pm$.
Hence
\begin{align*}
(\tecxt\text{b}) &= 0,
\end{align*}
because $a_1$ is applied to a state, that does not contain a 1.
The others are
\begin{align*}
a_4^\dagger a_7^\dagger a_5 a_1 a_1^\dagger a_3^\dagger a_5^\dagger a_8^\dagger\Ket | {0} \rangle &= \sqrt{(4-2+2)!}\Ket | {4,7,3,8} \rangle\\
a_5^\dagger a_3^\dagger a_8^\dagger a_1 a_3^\dagger a_2^\dagger a_1^\dagger a_7^\dagger\Ket | {0} \rangle &=\sqrt{(8-1)!}\Ket | {5,3,8,3,2,7} \rangle
\end{align*}
\section{Number operator}
We know that $\klam\left[ {a_n,a_m^\dagger} \right]=\delta_{nm}$ for bosons and $\klammer\left\{{c_n,c_m^\dagger} \right\}=\delta_{nm}$ for fermions.
Then we can show the following relations for the number operator $\hat{N}:=\sum_j a_j^\dagger a_j$, by making use of $\klam\left[ {A,BC} \right]=B\klam\left[ {A,C} \right]+\klam\left[ {A,B} \right]C$ and $\klam\left[ {A,BC} \right]=B\klammer\left\{{A,C} \right\}-\klammer\left\{{A,B} \right\}C$:
\begin{align*}
\klam\left[ {\hat{N},a_k^\dagger} \right]&=\sum_j\klam\left[ {a_j^\dagger a_j,a_k^\dagger} \right]=\sum_j a_j^\dagger \klam\left[ {a_j,a_k^\dagger} \right]=a_k^\dagger\\
\klam\left[ {\hat{N},a_k} \right]&=\sum_j\klam\left[ {a_j^\dagger a_j,a_k} \right]=\sum_j \klam\left[ {a_j^\dagger,a_k} \right]a_j =-a_k\\
\klam\left[ {\hat{N},c_k^\dagger} \right]&=\sum_j\klam\left[ {c_j^\dagger c_j,c_k^\dagger} \right]=\sum_jc_j^\dagger\klammer\left\{{c_j,c_k^\dagger} \right\}=c_k^\dagger\\
\klam\left[ {\hat{N},c_k} \right]&=\sum_j\klam\left[ {c_j^\dagger c_j,c_k} \right]=\sum_j-\klammer\left\{{c_j^\dagger,c_k} \right\}c_j=-c_k
\end{align*}
Here we used $\klam\left[ {a_n^\dagger,a_m^\dagger} \right]=0=\klammer\left\{{c_n^\dagger,c_m^\dagger} \right\}$ and $\klam\left[ {a_n,a_m} \right]=0=\klammer\left\{{c_n,c_m} \right\}$.
\\
Let the general Hamiltonian be
\begin{align*}
H&=\sum_{ij}T_{ij}a_i^\dagger a_j+\frac{1}{2}\sum_{ijkl}V_{ijkl}a_i^\dagger a_j^\dagger a_l a_k.
\end{align*}
Then this commutes with the number operator. 
Consider first the kinetic part, i.e.
\begin{align*}
\klam\left[ {\hat{N},T} \right]&=\sum_{ij}T_{ij}\klam\left[ {\hat{N},a_i^\dagger a_j} \right]=\sum_{ij}T_{ij}a_i^\dagger (-a_j)+a_i^\dagger a_j=0.
\end{align*}
The potential term is more complicated, so we will only write down the relevant part
\begin{align*}
\klam\left[ {\hat{N},a_i^\dagger a_j^\dagger a_l a_k} \right]&=a_i^\dagger a_j^\dagger \klam\left[ {\hat{N},a_l a_k} \right]+\klam\left[ {\hat{N},a_i^\dagger a_j^\dagger} \right]a_l a_k=\\
&=a_i^\dagger a_j^\dagger\kla\left( {a_l(-a_k) + (-a_l)a_k} \right)+\kla\left( {a_i^\dagger (+a_j^\dagger) + (+a_i^\dagger) a_j^\dagger} \right)a_l a_k=0.
\end{align*}
\section{Second quantization I}
Let the discrete basis be given by $\Ket | {\psi_k} \rangle$.
Then we can express the operator as
\begin{align*}
\mathcal{O}&=\sum_i\mathcal{O}_i=\sum_{jk}O_{jk}\sum_i a_j^\dagger a_k,
\end{align*}
where the $O_{jk}=\braket\langle {\psi_j} | {\mathcal{O}} | {\psi_k}\rangle$ are the matrix elements. 
FIXME: how to simplify these further?
\\
For the momentum operator we have
\begin{align*}
P_{p'p}&=\braket\langle {p'} | {P} | {p}\rangle=p'\delta (p'-p)
\end{align*}
and therefore
\begin{align*}
P&=\int_{p}\int_{p'}P_{p'p}a_{p'}^\dagger a_p=\intr\int_{\mathbb{R}^{3}}{pa_{p}^\dagger a_p}\,\mathrm{d}^{3}p.
\end{align*}
\section{Second quantization II}
4\texttt{D}\textsc{D}\textsf{D}\textsl{D}\textit{D}D
\section{$N$ bosons in a cube of edge length $L$}
\begin{align*}
k_n&=\begin{cases}
+1&\tecxt\text{if}\ \delta_{nm}>1/2 \\
-1&\tecxt\text{if}\ \delta_{nm}<-1/2 \\
\ 0 &\tecxt\text{else}
\end{cases} .\numberthis\label{eqn:RulesForKN}
\end{align*}
\end{document}